% BHU PYQ -- Manually transcribed & error-corrected
% Paper: GRA/GRB-605: Regional Development and Planning
% Exam: B.A./B.Sc. (Hons.) Semester VI Examination, 2013-14

\documentclass[11pt,a4paper]{article}
\usepackage[a4paper,top=2.5cm,bottom=2.5cm,left=2.8cm,right=2.8cm,headheight=14pt]{geometry}
\usepackage{amsmath,amssymb}
\usepackage[T1]{fontenc}
\usepackage[utf8]{inputenc}
\usepackage{lmodern,microtype,fancyhdr,enumitem,hyperref,lastpage,parskip}

\hypersetup{colorlinks=true,urlcolor=black,linkcolor=black,
  pdftitle={GRA/GRB-605 Regional Development and Planning -- BHU B.A./B.Sc. Sem VI 2013-14}}
\pagestyle{fancy}\fancyhf{}
\renewcommand{\headrulewidth}{0.4pt}\renewcommand{\footrulewidth}{0.4pt}
\fancyhead[L]{\small   \textit{Banaras Hindu University}}
\fancyhead[R]{\small   \textit{GRA/GRB-605: Regional Development and Planning}}
\fancyfoot[C]{\small Page  \thepage\ of \pageref{LastPage}}


\newlist{parts}{enumerate}{2}
\setlist[parts,1]{label=(\alph*),leftmargin=*,itemsep=5pt,topsep=3pt}

\newcommand{\pts}[1]{\hfill   \textit{[#1]}}


\begin{document}

\begin{center}
  {\large\bfseries BANARAS HINDU UNIVERSITY}\\[2pt]
  {\normalsize B.A./B.Sc. (Hons.) --- Semester VI Examination, 2013-14}\\[6pt]
  \rule{0.8\linewidth}{0.5pt}\\[6pt]
  {\large\bfseries Paper: GRA/GRB-605 --- Regional Development and Planning (Elective)}\\[4pt]
  {\normalsize Time: 3 Hours \hfill Maximum Marks: 70}\\[2pt]
  {\small Ref. No.: BS/Sem VI/290}
\end{center}
\rule{\linewidth}{0.4pt}
\smallskip



\noindent   \textit{Instructions: Answer   \textbf{five} questions in all. Question no. 1, carrying 20 marks, is compulsory. Remaining questions are of equal marks.}
\bigskip




\noindent   \textbf{Question 1.} Write short answer of the following: \pts{20}

\begin{parts}
  \item Local-level planning
  \item Participatory planning
  \item Planning region
  \item Panchayat Raj
  \item National capital region
\end{parts}

\medskip


\noindent   \textbf{Question 2.} Discuss the meaning and scope of regional development and planning. \pts{12.5}

\medskip


\noindent   \textbf{Question 3.} Examine the theory of regional development of either Myrdal or Perroux. \pts{12.5}

\medskip


\noindent   \textbf{Question 4.} Describe various approaches to regional development. \pts{12.5}

\medskip


\noindent   \textbf{Question 5.} Explain the evolution of regional Planning in India with examples. \pts{12.5}

\medskip


\noindent   \textbf{Question 6.} What is multi-level planning? Discuss its significance for development planning in India. \pts{12.5}

\medskip


\noindent   \textbf{Question 7.} "Infrastructure plays key role in regional development". Explain with suitable examples. \pts{12.5}

\medskip


\noindent   \textbf{Question 8.} Discuss the patterns of regional development and imbalances in India and identify the factors responsible for them. \pts{12.5}

\medskip


\noindent   \textbf{Question 9.} Critically examine the development planning of either Eastern Uttar Pradesh or North-East India. \pts{12.5}

\vfill\rule{\linewidth}{0.4pt}

\begin{center}\small
\href{https://bhu-exam.vercel.app/}{Click here to visit our website}\\[4pt]
\href{https://me-aryan.vercel.app/}{Click here to visit our contributor}\\[4pt]
\href{https://quashan-me.vercel.app/}{Click here to visit our contributor}
\end{center}
\end{document}
